\chapter{Как чинить машину без схемы}
% полная версия: chapters/17_debugging.tex

\pencilfig{17}{\draw[pen=pcGreen, wash=pcGreen!60] (0,0) rectangle (4,2);
\foreach \x/\y in {0.4/0.4,1.4/1.2,2.6/0.5} \draw[pen, wash=pcGray!70] (\x,\y) rectangle (\x+0.6,\y+0.4);
\draw[pen] (1.0,0.6) -- (1.4,0.6) -- (1.4,1.2); \draw[pen] (2.0,1.4) -- (2.9,1.4) -- (2.9,0.9); \draw[pen] (0.7,0.8) -- (0.7,1.6) -- (1.4,1.6);
\draw[pen=pcRed, line width=1.2pt] (5.2,2.7) -- (4.1,2.1); \draw[pen=pcRed] (4.1,2.1) -- (2.95,1.42);}{Плата без схемы: щуп слышит лишь несколько проводов из миллиардов.}

Представьте, что вам дали работающую плату без схемы и попросили её починить. У инженера есть четыре инструмента: щуп, чтобы послушать сигнал; генератор, чтобы подать свой сигнал; возможность отключить деталь и посмотреть, что сломается; и доступ к регистрам управления. Ровно этими инструментами изучают и мозг.

\section*{Щуп на тысячу нейронов}

Электрод, воткнутый в мозг, слышит щелчки нескольких нейронов рядом. Казалось бы, из восьмидесяти шести миллиардов тысяча электродов --- ничто. Но их хватает, чтобы управлять курсором на экране. Почему? Соседние нейроны шумят похоже, и активность, управляющая рукой, на самом деле описывается всего десятком-двумя общих переменных. Не нужно слушать всех, достаточно услышать эти несколько «голосов».

Сырой сигнал с тысячи электродов --- это сотни миллионов бит в секунду: радиоканалу из-под черепа не пропустить. А сами щелчки несут примерно в тысячу раз меньше. Поэтому щелчки выгодно выделять прямо на имплантате и отправлять наружу только их --- тот же приём сжатия, что в глазу.

\section*{Что делает Neuralink}

Компания Neuralink взялась за инженерную сторону задачи: тонкие гибкие нити вместо жёстких иголок, робот, который вживляет их, и герметичный беспроводной имплантат. В её опубликованной работе описана система, соединённая с компьютером кабелем, а сжатие и беспроводная связь названы планами. Данные о человеке с парализованными руками, который управляет курсором, пока известны в основном по отчётам самой компании; по ним в имплантате около тысячи каналов на нескольких десятках нитей, и часть нитей после вживления сместилась. Сама идея не нова: академические группы с жёсткими решётками примерно из сотни электродов уже давали людям управлять курсором, «писать» буквы силой мысли и превращать попытки говорить в текст на экране со скоростью около шестидесяти слов в минуту.

\section*{Чтение и запись}

Декодер читает частоты щелчков многих нейронов и получает меньше десяти бит в секунду --- ничтожную долю того, что несут сами нейроны. Зато и мозг, и декодер учатся друг у друга: человек приспосабливается к декодеру, декодер --- к человеку. Два ученика, подстраивающихся друг под друга, --- тонкая задача: даже когда каждый в отдельности устойчив, вместе они могут раскачаться.

Записывать в мозг труднее, чем читать. Ток через электрод возбуждает всё вокруг, а не один нужный нейрон. Можно «нарисовать» точки света, как светящиеся пиксели, и человек различит простые буквы. Но записать сжатый код глаза так, чтобы мозг увидел картинку, пока не умеет никто.

\section*{Чего щуп не видит}

Электрод слышит телефонную сеть данных. Но радиостанции --- дофамин, серотонин, норадреналин, ацетилхолин --- это концентрации веществ, их щуп не слышит. Не видит он и силы связей, где хранится вся память. И не видит боли так, как её чувствует человек. Отладчик, который видит данные, но не видит регистров, всегда будет удивляться, почему один и тот же вход даёт разные ответы.

\fullversion{17}
